\scanpage{139}
\mspage{V.11}
% archive sheet number: 138
% manuscript lines 1-3
est surjectif si $i\geq\dim Y_{\unsure{\bar\eta|\bar y}}=n-d-1$, \deleted{\unsure{bi|ou}} donc si $i+2(d+1)\geq n+d+1$, \inserted{et bijectif si $i>n-d-1$ i.e. si $i+2(d+1)>n+d+1$,}
% manuscript lines 4-7: cancelled argument
\deleted{Donc le fait que $I$ opère trivialement sur $H^i(X_{\bar\eta},\Lambda_{\unsure{\nu|v}})$ est acquis puisque le fait qu'il opère trivialement sur $H^i(Y_{\bar\eta},\Lambda_\nu)$, qui est \unsure{une} conséquence bien connue du fait que $Y$ est propre et lisse sur $S$ (SGA 4 XVI ---).}
\inserted{Le diagramme commutatif}
\inserted{\diagram{\begin{tikzcd}[ampersand replacement=\&]
H^i(\deleted{\unsure{X}}Y_{\bar s},\Lambda)\arrow[r]\arrow[d,"\varphi_Y"']\&H^{i+2(d+1)}(X_0)\arrow[d,"\varphi_X"]\\
H^i(Y_{\bar\eta},\Lambda)\arrow[r]\&H^{i+2(d+1)}(\deleted{\unsure{Y}}X_{\bar\eta})
\end{tikzcd}}}
% slanted
\margin{donne alors le résultat voulu ($\varphi_Y$ est bijectif par SGA 4 \uline{XVI} ---)}

% manuscript lines 8-12
\deleted{\unsure{2.}}\uline{Remarque} 2.8. L'assertion 2.3 suggère la \deleted{validité de} la question suivante. Soit $X$ \inserted{$=\operatorname{Spec}A$} un schéma strictement local \inserted{de dim. \deleted{\unsure{un}} $n$}, $t\in\unsure{\mathfrak m|\mathfrak r}_A$, \deleted{un élément de et} $U=X-V(t)=\operatorname{Spec}A_t$, et soit pour les entiers $m$ premiers à la car. résiduelle
\[
\unsure{X_m|X_{(m)}}=\operatorname{Spec}A[T]/(T^m-t);\qquad U(m)=X_m\times_X U,
\]
enfin soit
\[
\deleted{\unsure{X}}U(\infty)=\varprojlim U(m).
\]
% manuscript lines 13-18
Est-il vrai que la dim. cohomologique de $U(\infty)$, pour des coefficients premiers à \deleted{la} $p$, est $\leq n-1$ ? [La conjecture standard SGA 4 XIV \unsure{et|e} \deleted{donnerait par} donnerait \inserted{pour $\ell\ne p$ :} $\operatorname{cd}_\ell(U(m))\leq\unsure{n}$ pour $\forall\,m$, donc $\operatorname{cd}_\ell(U(\infty))\leq n$, et non $n-1$ !
